% Encuesta de evaluación experta – A tu Ritmo
% Paquetes necesarios en el preámbulo:
%   \usepackage[spanish]{babel}
%   \usepackage{booktabs}
%   \usepackage{tabularx}
%   \usepackage{array}
%   \usepackage{multirow}
% Si el documento no los tiene, descomente el bloque de abajo para compilar este archivo solo.

% \documentclass[11pt]{article}
% \usepackage[utf8]{inputenc}
% \usepackage[spanish]{babel}
% \usepackage[margin=2.2cm]{geometry}
% \usepackage{booktabs,tabularx,array,multirow}
% \begin{document}

\newcolumntype{C}{>{\centering\arraybackslash}p{0.55cm}}
\renewcommand{\arraystretch}{1.35}

% ---------------------------------------------------------------
\begin{table}[htbp]
\centering
\caption{Escala de valoración utilizada en los bloques A, B y C}
\label{tab:escala}
\begin{tabular}{ccccc}
\toprule
\textbf{1} & \textbf{2} & \textbf{3} & \textbf{4} & \textbf{5} \\
Totalmente en desacuerdo & En desacuerdo & Neutral & De acuerdo & Totalmente de acuerdo \\
\bottomrule
\end{tabular}
\end{table}

% ---------------------------------------------------------------
\begin{table}[htbp]
\centering
\caption{Datos del evaluador}
\label{tab:perfil}
\begin{tabularx}{\textwidth}{lX}
\toprule
\textbf{Campo} & \textbf{Respuesta} \\
\midrule
Nombre completo & \\
Institución & \\
Área de especialidad & Educación especial / Musicoterapia / Educación musical / Terapia ocupacional / Fonoaudiología / Psicología / Desarrollo de software o UX / Otra \\
Años de experiencia & \\
Fecha de evaluación & \\
Sesiones observadas & \\
Niños o niñas observados & \\
Dispositivo usado & Tablet Android / Teléfono Android / Computador Windows \\
Actividades observadas & Simón (memoria) / Sigue el ritmo / Tempo Tap / Banda musical \\
\bottomrule
\end{tabularx}
\end{table}

% ---------------------------------------------------------------
\begin{table}[htbp]
\centering
\caption{Bloque A. Evaluación de la aplicación (Escala de Usabilidad del Sistema, SUS)}
\label{tab:bloqueA}
\begin{tabularx}{\textwidth}{lXCCCCC}
\toprule
\textbf{N.\textsuperscript{o}} & \textbf{Afirmación} & \textbf{1} & \textbf{2} & \textbf{3} & \textbf{4} & \textbf{5} \\
\midrule
A1  & Creo que a los niños y niñas les gustaría usar esta aplicación con frecuencia. & & & & & \\
A2  & Encontré la aplicación innecesariamente compleja. & & & & & \\
A3  & La aplicación fue fácil de usar. & & & & & \\
A4  & Se necesita el apoyo de un técnico para poder usar esta aplicación. & & & & & \\
A5  & Las funciones de la aplicación (perfiles, menú y juegos) están bien integradas. & & & & & \\
A6  & Hay demasiada inconsistencia en la aplicación. & & & & & \\
A7  & La mayoría de los niños y niñas aprenderían a usar esta aplicación muy rápidamente. & & & & & \\
A8  & La aplicación es muy difícil de usar. & & & & & \\
A9  & Los niños y niñas se mostraron seguros y confiados al usar la aplicación. & & & & & \\
A10 & Fue necesario aprender muchas cosas antes de poder usar la aplicación. & & & & & \\
\midrule
\multicolumn{7}{p{0.95\textwidth}}{\textbf{Observaciones sobre la interfaz} (botones, colores, sonidos, retroalimentación):} \\[1.6cm]
\bottomrule
\end{tabularx}
\end{table}

% ---------------------------------------------------------------
\begin{table}[htbp]
\centering
\caption{Bloque B. Actividades rítmicas y valor educativo}
\label{tab:bloqueB}
\begin{tabularx}{\textwidth}{lXCCCCC}
\toprule
\textbf{N.\textsuperscript{o}} & \textbf{Afirmación} & \textbf{1} & \textbf{2} & \textbf{3} & \textbf{4} & \textbf{5} \\
\midrule
\multicolumn{7}{l}{\textit{Diseño de las actividades}} \\
B1  & Las actividades rítmicas son adecuadas a la edad y al nivel de desarrollo de los niños y niñas con síndrome de Down. & & & & & \\
B2  & Las instrucciones de cada juego se comprenden sin explicación adicional del adulto. & & & & & \\
B3  & El tempo y la velocidad de los patrones son apropiados para los usuarios. & & & & & \\
B4  & La dificultad aumenta de forma gradual y a un ritmo adecuado. & & & & & \\
B5  & La retroalimentación visual y sonora ayuda a corregir el ritmo. & & & & & \\
\midrule
\multicolumn{7}{l}{\textit{Valor educativo}} \\
B6  & Las actividades trabajan habilidades rítmicas concretas: pulso, acento e imitación de patrones. & & & & & \\
B7  & Los juegos favorecen la atención y la memoria secuencial. & & & & & \\
B8  & Las actividades estimulan la coordinación motora (toques y gestos al compás). & & & & & \\
B9  & Los contenidos pueden integrarse a la planificación de aula o de terapia. & & & & & \\
B10 & Las actividades mantienen la motivación durante toda la sesión. & & & & & \\
\midrule
\multicolumn{7}{p{0.95\textwidth}}{\textbf{Comentarios sobre las actividades} (qué juego aporta más, qué cambiaría):} \\[1.6cm]
\bottomrule
\end{tabularx}
\end{table}

% ---------------------------------------------------------------
\begin{table}[htbp]
\centering
\caption{Bloque C. Progreso y aprendizaje observado en el ritmo}
\label{tab:bloqueC}
\begin{tabularx}{\textwidth}{lXCCCCC}
\toprule
\textbf{N.\textsuperscript{o}} & \textbf{Afirmación} & \textbf{1} & \textbf{2} & \textbf{3} & \textbf{4} & \textbf{5} \\
\midrule
C1 & Mejoraron su capacidad para mantener un pulso estable. & & & & & \\
C2 & Reprodujeron patrones rítmicos con mayor precisión que al inicio. & & & & & \\
C3 & Lograron sincronizar sus toques con la música o el metrónomo. & & & & & \\
C4 & Recordaron secuencias más largas a medida que avanzaban. & & & & & \\
C5 & Aumentó el tiempo de atención sostenida durante la actividad. & & & & & \\
C6 & Ganaron autonomía: necesitaron menos ayuda del adulto para jugar. & & & & & \\
C7 & Mostraron disfrute e interés por volver a jugar. & & & & & \\
C8 & Los reportes de partidas reflejan el progreso que observé en las sesiones. & & & & & \\
\midrule
\multicolumn{7}{p{0.95\textwidth}}{\textbf{Ejemplo concreto de progreso o dificultad observada:}} \\[1.6cm]
\bottomrule
\end{tabularx}
\end{table}

% ---------------------------------------------------------------
\begin{table}[htbp]
\centering
\caption{Dictamen del experto}
\label{tab:dictamen}
\begin{tabularx}{\textwidth}{XC}
\toprule
\multicolumn{2}{l}{\textbf{¿Recomienda el uso de A tu Ritmo con niños y niñas con síndrome de Down?}} \\
\midrule
Sí, tal como está. Puede usarse en aula o terapia sin cambios. & \\
Sí, con ajustes menores. Requiere cambios puntuales antes de usarse. & \\
Requiere ajustes importantes. Debe revisarse antes de aplicarse. & \\
No la recomiendo para esta población. & \\
\midrule
\multicolumn{2}{p{0.95\textwidth}}{\textbf{Principales fortalezas:}} \\[1.6cm]
\multicolumn{2}{p{0.95\textwidth}}{\textbf{Mejoras sugeridas:}} \\[1.6cm]
\bottomrule
\end{tabularx}
\end{table}

% ---------------------------------------------------------------
\begin{table}[htbp]
\centering
\caption{Cálculo de puntajes por bloque}
\label{tab:puntajes}
\begin{tabularx}{\textwidth}{lXl}
\toprule
\textbf{Bloque} & \textbf{Cálculo} & \textbf{Rango} \\
\midrule
A (SUS) & Ítems impares: $x_i - 1$. Ítems pares: $5 - x_i$. Se suman los 10 valores y se multiplica por $2{,}5$. & 0--100 \\
B y C & Promedio de $\dfrac{x_i - 1}{4} \times 100$ sobre los ítems del bloque. & 0--100 \\
\midrule
\multicolumn{3}{p{0.95\textwidth}}{\textit{Interpretación del SUS:} $\geq 85$ excelente; 72--84 buena; 68--71 aceptable; 51--67 marginal; $< 51$ deficiente. \textit{Bloques B y C:} $\geq 80$ muy alto; 60--79 alto; 40--59 medio; $< 40$ bajo.} \\
\bottomrule
\end{tabularx}
\end{table}

% \end{document}
